\documentclass[11pt,a4paper]{article}

\usepackage[T1]{fontenc}
\usepackage[utf8]{inputenc}
\usepackage{libertinus}
\usepackage[final]{microtype}
\usepackage[a4paper,margin=26mm,headheight=15pt]{geometry}
\usepackage{amsmath,amssymb,amsthm,mathtools}
\usepackage{bm,mathrsfs}
\usepackage{booktabs,longtable,array,tabularx}
\usepackage{enumitem}
\usepackage{xcolor}
\usepackage[most]{tcolorbox}
\usepackage{listings}
\usepackage{fancyhdr}
\usepackage{xurl}
\usepackage{hyperref}
\usepackage[nameinlink,capitalise,noabbrev]{cleveref}

\definecolor{deepolive}{RGB}{67,79,45}
\definecolor{softolive}{RGB}{236,239,226}
\definecolor{darkteal}{RGB}{31,83,84}
\definecolor{warmgray}{RGB}{92,87,78}
\definecolor{softcoral}{RGB}{246,225,219}
\definecolor{softgray}{RGB}{247,247,247}
\definecolor{linkblue}{RGB}{31,76,122}

\hypersetup{
  colorlinks=true,
  linkcolor=deepolive,
  citecolor=darkteal,
  urlcolor=linkblue,
  pdftitle={The Full Asymptotic Transseries of the Swing Factorial and Its Inverse},
  pdfauthor={Research note prepared for the ProveIt project},
  pdfsubject={Swinging factorial, gamma ratios, Bernoulli and Bell coefficients, Borel summation, Lambert W, asymptotic inversion}
}

\pagestyle{fancy}
\fancyhf{}
\fancyhead[L]{\small Swing factorial and its inverse}
\fancyhead[R]{\small Asymptotic transseries}
\fancyfoot[C]{\thepage}
\renewcommand{\headrulewidth}{0.4pt}

\setlength{\parindent}{0pt}
\setlength{\parskip}{0.52em}
\setlist{itemsep=2pt,topsep=4pt,parsep=1pt,leftmargin=2em}
\allowdisplaybreaks[3]
\numberwithin{equation}{section}
\emergencystretch=2em

\newtheorem{theorem}{Theorem}[section]
\newtheorem{lemma}[theorem]{Lemma}
\newtheorem{proposition}[theorem]{Proposition}
\newtheorem{corollary}[theorem]{Corollary}
\newtheorem{identity}[theorem]{Identity}
\theoremstyle{definition}
\newtheorem{definition}[theorem]{Definition}
\newtheorem{example}[theorem]{Example}
\newtheorem{algorithm}[theorem]{Algorithm}
\theoremstyle{remark}
\newtheorem{remark}[theorem]{Remark}

\newtcolorbox{mainresultbox}[1][]{
  enhanced,breakable,colback=softolive,colframe=deepolive,
  boxrule=0.8pt,arc=1.2mm,left=3mm,right=3mm,top=2.5mm,bottom=2.5mm,
  title={#1},fonttitle=\bfseries
}
\newtcolorbox{methodbox}[1][]{
  enhanced,breakable,colback=softgray,colframe=warmgray,
  boxrule=0.5pt,arc=1mm,left=3mm,right=3mm,top=2mm,bottom=2mm,#1
}
\newtcolorbox{cautionbox}[1][]{
  enhanced,breakable,colback=softcoral,colframe=warmgray,
  boxrule=0.6pt,arc=1mm,left=3mm,right=3mm,top=2mm,bottom=2mm,#1
}

\lstdefinestyle{pythonstyle}{
  language=Python,
  basicstyle=\ttfamily\footnotesize,
  columns=fullflexible,
  keepspaces=true,
  frame=single,
  framerule=0.4pt,
  breaklines=true,
  showstringspaces=false,
  keywordstyle=\bfseries,
  commentstyle=\itshape,
  xleftmargin=0.4em,
  xrightmargin=0.4em
}

\newcommand{\e}{\mathrm e}
\newcommand{\dd}{\,\mathrm d}
\newcommand{\ii}{\mathrm i}
\newcommand{\Ord}{\mathcal O}
\newcommand{\littleo}{\mathrm o}
\newcommand{\coeff}[2]{[\,#1^{#2}\,]}
\newcommand{\Coeff}[1]{[\,#1\,]}
\newcommand{\Borel}{\mathcal B}
\newcommand{\Laplace}{\mathcal L}
\newcommand{\SF}{\operatorname{sf}}
\newcommand{\Disc}{\operatorname{Disc}}
\newcommand{\W}{W}
\newcommand{\Log}{\operatorname{Log}}
\newcommand{\Res}{\operatorname{Res}}
\newcommand{\BellY}{Y}
\newcommand{\StirlingFirst}[2]{\genfrac{[}{]}{0pt}{}{#1}{#2}}

\title{\textbf{The Full Asymptotic Transseries of the Swing Factorial and Its Inverse}\\[0.35em]
\Large Gamma-Ratio Normal Forms, Bernoulli--Bell Coefficient Calculus,\\
Borel Geometry, and Lambert-Normalized Reversion}
\author{Research note prepared for the ProveIt project}
\date{3 September 2026}

\begin{document}
\maketitle

\begin{abstract}
The swinging factorial (or swing factorial)
\[
  \SF(n)=\frac{n!}{\lfloor n/2\rfloor!^2}
\]
is not governed by a single smooth asymptotic scale: its even and odd
subsequences have different algebraic prefactors, and the complete sequence
falls sharply after every sufficiently large odd term.  We therefore treat
parity as a genuine two-component transseries sector.  With
\(\sigma=-1\) for the even phase and \(\sigma=+1\) for the odd phase, the
canonical analytic branches are
\[
 \mathcal S_\sigma(x)
 =\frac{\Gamma(x+1)}{\Gamma\!\left(x/2+(3-\sigma)/4\right)^2}
 =\frac{2^x}{\sqrt\pi}
   \left(\frac{\Gamma(x/2+1)}{\Gamma(x/2+1/2)}\right)^\sigma.
\]
They admit the exact normal form
\[
 \mathcal S_\sigma(x)
 =C_\sigma 2^x x^{\sigma/2}\exp\!\bigl(\sigma D(x)\bigr),
 \qquad C_\sigma=\frac{2^{-\sigma/2}}{\sqrt\pi},
\]
where
\[
 D(x)=\frac12\int_0^\infty \e^{-xt}\frac{\tanh(t/2)}{t}\,\dd t.
\]
Thus the Borel transform is the elementary meromorphic function
\(\widehat D(t)=\tanh(t/2)/(2t)\).  This gives, simultaneously, the
all-orders logarithmic expansion, an exact signed remainder, its optimal
truncation scale, and the complete Borel singularity lattice.

The logarithmic coefficients are
\[
 d_{2r+1}=\frac{(2^{2r+2}-1)B_{2r+2}}{(2r+1)(2r+2)},\qquad d_{2r}=0.
\]
Exponentiation gives every coefficient of the swing factorial itself by a
complete exponential Bell polynomial, an explicit partition sum, or a
triangular recurrence.  For inversion, the dominant equation
\(L=\lambda x+(\sigma/2)\log x\), \(\lambda=\log2\), is solved exactly by
Lambert's \(W\)-function.  The real large branches require \(W_0\) in the
odd sector but \(W_{-1}\) in the even sector.  Around that Lambert core the
inverse has the constant-coefficient expansion
\[
 \mathcal I_\sigma(y)
 \sim X_\sigma(y)+\sum_{n\ge1}\frac{c_n^{(\sigma)}}{X_\sigma(y)^n},
\]
and we give both a Lagrange--B\"urmann coefficient formula and an ordinary
Bell-polynomial recurrence for arbitrary \(c_n^{(\sigma)}\).  A second,
fully expanded power--logarithmic form
\[
 \mathcal I_\sigma(y)
 \sim\frac1\lambda\left(
 L_\sigma-\frac\sigma2H_\sigma+
 \sum_{n\ge1}\frac{P_n^{(\sigma)}(H_\sigma)}{L_\sigma^n}
 \right)
\]
is developed in parallel.  General formulae for the polynomials
\(P_n^{(\sigma)}\) are expressed through coefficient extraction and ordinary
Bell polynomials; their pure Lambert block is written explicitly with
unsigned Stirling numbers of the first kind.  We finish with Borel--Stokes
transport under inversion, rigorous residual error estimates, numerical
validation, and symbolic-generation code.
\end{abstract}

\begin{center}
\small\textbf{Keywords.}
Swinging factorial; central binomial coefficient; gamma ratio; Bernoulli
numbers; Bell polynomials; Stirling numbers; Borel summation; resurgent
asymptotics; Lambert \(W\); Lagrange--B\"urmann inversion; power--logarithmic
transseries.
\end{center}

\tableofcontents

\section{Introduction and statement of the main results}

The sequence OEIS A056040 begins
\[
 1,1,2,6,6,30,20,140,70,630,252,2772,\ldots
\]
and is commonly called the \emph{swinging factorial}.  Its elementary
formulae, generating functions, and arithmetic interpretations are recorded
in OEIS and in Luschny's work on swinging factorials
\cite{OEISA056040,LuschnySwing}.  The existing OEIS entry A219931 also
records a finite portion of the asymptotic expansion of its logarithm
\cite{OEISA219931}.  The purpose of this article is to put the problem into a
complete coefficient-calculus and transseries framework, and then to carry
out the substantially subtler inverse problem.

Two features dictate the organization.

\begin{enumerate}
\item The even and odd subsequences are respectively
\[
 \SF(2m)=\binom{2m}{m},\qquad
 \SF(2m+1)=(2m+1)\binom{2m}{m}.
\]
Their leading powers of the index are opposite: \(n^{-1/2}\) versus
\(n^{+1/2}\).
\item The full sequence is not eventually monotone.  Indeed
\[
 \frac{\SF(2m+1)}{\SF(2m)}=2m+1,
 \qquad
 \frac{\SF(2m+2)}{\SF(2m+1)}=\frac{2}{m+1}.
\]
Consequently there is no canonical one-valued inverse of the unsplit integer
sequence.
\end{enumerate}

The correct asymptotic object is therefore a pair of monotone analytic
branches.  Throughout,
\[
 \sigma=-1\quad\hbox{means even phase},\qquad
 \sigma=+1\quad\hbox{means odd phase},
 \qquad \lambda:=\log2.
\]
At an integer \(n\), its phase is
\[
 \sigma_n=(-1)^{n+1}=-(-1)^n.
\]
It is useful to regard calculations as taking place in the parity algebra
\[
 \mathscr R=\mathbb Q[\lambda,\lambda^{-1},\sigma]/(\sigma^2-1).
\]
This records both sectors in one formula without pretending that parity is a
small perturbation.

\begin{mainresultbox}[Forward transseries]
For every fixed \(N\ge1\), uniformly in closed subsectors of \(\Re x>0\),
\[
 \mathcal S_\sigma(x)
 =C_\sigma 2^x x^{\sigma/2}
 \left(\sum_{n=0}^{N-1}\frac{\sigma^n A_n}{x^n}
       +\Ord(x^{-N})\right),
 \qquad C_\sigma=\frac{2^{-\sigma/2}}{\sqrt\pi},
\]
where
\[
 A_n=\frac1{n!}\BellY_n(1!d_1,2!d_2,\ldots,n!d_n)
 =\sum_{\substack{k_1,k_2,\ldots\ge0\\
                   \sum_{j\ge1}jk_j=n}}
   \prod_{j\ge1}\frac{d_j^{k_j}}{k_j!},
\]
\[
 d_{2r+1}=\frac{(2^{2r+2}-1)B_{2r+2}}{(2r+1)(2r+2)},
 \qquad d_{2r}=0.
\]
Equivalently,
\[
 A_0=1,\qquad nA_n=\sum_{j=1}^n j d_jA_{n-j}.
\]
The logarithmic series is Borel summable to an exact elementary Laplace
integral, and its remainder after \(N\) nonzero terms has the sign of the
first omitted term and smaller magnitude.
\end{mainresultbox}

\begin{mainresultbox}[Lambert-normalized inverse transseries]
Put
\[
 L_\sigma(y)=\log\frac{y}{C_\sigma}
 =\log y+\frac12\log\pi+\frac\sigma2\log2.
\]
The exact inverse of the dominant block
\(L_\sigma=\lambda x+(\sigma/2)\log x\) is
\[
 X_{\sigma,k}(y)
 =\frac{\sigma}{2\lambda}
  W_k\!\left(2\sigma\lambda\e^{2\sigma L_\sigma(y)}\right).
\]
For the large positive real inverse,
\[
 X_+(y)=\frac{W_0(4\pi\lambda y^2)}{2\lambda},\qquad
 X_-(y)=-\frac{W_{-1}(-4\lambda/(\pi y^2))}{2\lambda}.
\]
Then
\[
 \mathcal I_\sigma(y)=\mathcal S_\sigma^{-1}(y)
 \sim X_\sigma(y)+\sum_{n\ge1}c_n^{(\sigma)}X_\sigma(y)^{-n}.
\]
Define
\[
 Q_\sigma(t)=\sigma\sum_{j\ge1}d_jt^j,
 \quad b=\frac\sigma2,
 \quad
 \Psi_\sigma(t,u)=-\frac1\lambda\left[
 b\log(1+tu)+Q_\sigma\!\left(\frac{t}{1+tu}\right)
 \right].
\]
The general coefficient is
\[
 c_n^{(\sigma)}
 =\sum_{m=1}^n\frac1m
   \Coeff{t^n u^{m-1}}\Psi_\sigma(t,u)^m,
\]
where the notation means ordinary coefficient extraction in both variables.
A completely triangular Bell-polynomial recurrence is given in
\cref{thm:lambert-centered}.
\end{mainresultbox}

\begin{mainresultbox}[Direct power--logarithmic inverse]
Let
\[
 H_\sigma(y)=\log\frac{L_\sigma(y)}{\lambda}.
\]
Then
\[
 \mathcal I_\sigma(y)
 \sim\frac1\lambda\left[
 L_\sigma-\frac\sigma2H_\sigma+
 \sum_{n\ge1}\frac{P_n^{(\sigma)}(H_\sigma)}{L_\sigma^n}
 \right],
\]
where \(P_n^{(\sigma)}\) is a polynomial of degree at most \(n\), generated
at arbitrary order by \cref{thm:direct-reversion}.  If the gamma-ratio
correction is suppressed, the remaining pure Lambert polynomial has the
closed form
\[
 P_n^{(0)}(H)=
 \frac{b^{n+1}}{n!}
 \sum_{m=1}^n(-1)^{m+1}(m-1)!
 \StirlingFirst{n}{m}\binom{n}{m-1}H^{n-m+1},
 \qquad b=\frac\sigma2.
\]
Thus the inverse coefficients sit exactly at the intersection of the
Bernoulli, Bell, and Stirling coefficient calculi.
\end{mainresultbox}

\section{Elementary structure and the two analytic branches}

\subsection{Equivalent definitions}

\begin{definition}[Swing factorial]
For \(n\in\mathbb N_0\), define
\begin{equation}\label{eq:swing-def}
 \SF(n):=\frac{n!}{\lfloor n/2\rfloor!^2}.
\end{equation}
\end{definition}

Writing \(n=2m\) or \(2m+1\) immediately gives
\begin{equation}\label{eq:evenodd-basic}
 \SF(2m)=\binom{2m}{m},\qquad
 \SF(2m+1)=(2m+1)\binom{2m}{m}.
\end{equation}
It also follows that
\begin{equation}\label{eq:one-step-rec}
 \SF(0)=1,\qquad
 \SF(n)=
 \begin{cases}
   n\,\SF(n-1),&n\text{ odd},\\[1mm]
   \dfrac4n\,\SF(n-1),&n\text{ even}.
 \end{cases}
\end{equation}
The product form in the OEIS definition is recovered by iterating
\cref{eq:one-step-rec}.

\subsection{Generating functions}

\begin{proposition}[Ordinary and exponential generating functions]
The swing factorial has exponential generating function
\begin{equation}\label{eq:egf}
 \sum_{n\ge0}\SF(n)\frac{z^n}{n!}=(1+z)I_0(2z),
\end{equation}
and ordinary generating function
\begin{equation}\label{eq:ogf}
 \sum_{n\ge0}\SF(n)z^n
 =\frac1{\sqrt{1-4z^2}}+\frac{z}{(1-4z^2)^{3/2}}
 =\frac{1+z/(1-4z^2)}{\sqrt{1-4z^2}}.
\end{equation}
\end{proposition}

\begin{proof}
From \cref{eq:evenodd-basic},
\[
 \sum_{m\ge0}\SF(2m)\frac{z^{2m}}{(2m)!}
 =\sum_{m\ge0}\frac{z^{2m}}{(m!)^2}=I_0(2z),
\]
and the odd exponential series is \(zI_0(2z)\).  For the ordinary series,
use
\[
 F(u):=\sum_{m\ge0}\binom{2m}{m}u^m=(1-4u)^{-1/2}.
\]
The odd part is
\[
 z\sum_{m\ge0}(2m+1)\binom{2m}{m}z^{2m}
 =z\bigl(2uF'(u)+F(u)\bigr)\big|_{u=z^2}
 =\frac{z}{(1-4z^2)^{3/2}}.
\]
\end{proof}

The two singularities \(z=\pm1/2\) in \cref{eq:ogf} have the same modulus.
That co-dominance is the generating-function origin of the parity sector:
no single nonoscillatory singular expansion can represent the full sequence.
Even/odd projection separates the two contributions exactly.

\subsection{Canonical gamma continuations}

\begin{definition}[Parity-resolved analytic branches]\label{def:branches}
For \(\sigma\in\{-1,+1\}\), define
\begin{equation}\label{eq:branch-gamma}
 \mathcal S_\sigma(x)
 :=\frac{\Gamma(x+1)}
 {\Gamma\!\left(x/2+(3-\sigma)/4\right)^2}.
\end{equation}
Thus \(\mathcal S_-(2m)=\SF(2m)\) and
\(\mathcal S_+(2m+1)=\SF(2m+1)\).
\end{definition}

The duplication formula gives the more symmetric identity
\begin{equation}\label{eq:branch-ratio}
 \mathcal S_\sigma(x)
 =\frac{2^x}{\sqrt\pi}
 \left(\frac{\Gamma(x/2+1)}{\Gamma(x/2+1/2)}\right)^\sigma.
\end{equation}
It follows directly from \cref{eq:branch-gamma} that
\begin{equation}\label{eq:step-two}
 \frac{\mathcal S_\sigma(x+2)}{\mathcal S_\sigma(x)}
 =\frac{16(x+1)(x+2)}{(2x+3-\sigma)^2}
 =\begin{cases}
   \dfrac{4(x+1)}{x+2},&\sigma=-1,\\[2mm]
   \dfrac{4(x+2)}{x+1},&\sigma=+1.
  \end{cases}
\end{equation}
The recurrence makes the common exponential scale \(4\) per two steps
transparent.

\begin{proposition}[Monotonicity and real inverses]\label{prop:monotone}
Both branches are strictly increasing for \(x>0\).  More precisely,
\begin{equation}\label{eq:logder}
 \frac{\dd}{\dd x}\log\mathcal S_\sigma(x)
 =\lambda+\frac\sigma2
 \left[\psi(x/2+1)-\psi(x/2+1/2)\right].
\end{equation}
The odd branch is strictly increasing on \([0,\infty)\); the even branch has
zero derivative at \(0\) and positive derivative thereafter.  Consequently
each has a unique large positive inverse, denoted \(\mathcal I_\sigma\).
\end{proposition}

\begin{proof}
The digamma difference in brackets is positive because \(\psi\) is
increasing.  Its derivative is
\(\tfrac12[\psi_1(x/2+1)-\psi_1(x/2+1/2)]<0\), since the trigamma function is
positive and decreasing.  At \(x=0\) the difference equals \(2\log2\).
Thus the derivative in \cref{eq:logder} is positive for \(\sigma=+1\), and
for \(\sigma=-1\) it is zero only at \(x=0\) and positive for \(x>0\).
\end{proof}

\begin{cautionbox}
A global interpolation such as
\[
 \widetilde{\mathcal S}(x)=
 \frac{1+\cos\pi x}{2}\mathcal S_-(x)
 +\frac{1-\cos\pi x}{2}\mathcal S_+(x)
\]
matches every integer value, but is highly oscillatory and has no canonical
one-valued large inverse.  The discontinuous interpolation obtained by
retaining \(\lfloor x/2\rfloor\) is no better.  The phase-resolved pair
\((\mathcal S_-,\mathcal S_+)\) is therefore not merely convenient; it is
the asymptotically stable notion of continuation and inversion.
\end{cautionbox}

\section{Exact logarithmic normal form}

\subsection{Separating the dominant scale}

Set
\begin{equation}\label{eq:D-def}
 D(x):=
 \log\frac{\Gamma(x/2+1)}{\Gamma(x/2+1/2)}
 -\frac12\log\frac{x}{2}.
\end{equation}
Then \cref{eq:branch-ratio} becomes the exact factorization
\begin{equation}\label{eq:exact-normal}
 \boxed{\quad
 \mathcal S_\sigma(x)=
 C_\sigma\e^{\lambda x}x^{\sigma/2}\e^{\sigma D(x)},
 \qquad
 C_\sigma=\frac{2^{-\sigma/2}}{\sqrt\pi}.
 \quad}
\end{equation}
At integer arguments this gives the exact parity form
\begin{equation}\label{eq:integer-log-exact}
 \log\SF(n)=n\lambda-\frac12\log\pi
 +\sigma_n\left(\frac12\log\frac n2+D(n)\right),
 \qquad \sigma_n=(-1)^{n+1}.
\end{equation}
The periodic multiplier is part of the leading transseries data, not an error
term.

\subsection{An exact Laplace representation}

\begin{theorem}[Elementary Borel--Laplace representation]\label{thm:laplace-D}
For \(\Re x>0\),
\begin{equation}\label{eq:D-laplace}
 D(x)=\frac12\int_0^\infty
       \e^{-xt}\frac{\tanh(t/2)}{t}\,\dd t
 =\int_0^\infty \e^{-xt}\widehat D(t)\,\dd t,
 \qquad
 \widehat D(t)=\frac{\tanh(t/2)}{2t}.
\end{equation}
The apparent singularity of \(\widehat D\) at \(t=0\) is removable, with
\(\widehat D(0)=1/4\).
\end{theorem}

\begin{proof}
Let
\[
 h(z)=\log\Gamma(z+1)-\log\Gamma(z+1/2)-\frac12\log z,
 \qquad z=x/2.
\]
Using the standard integral representation for a digamma difference and the
identity \(1/(2z)=\int_0^\infty \e^{-2zu}\,\dd u\), one obtains
\begin{align*}
 h'(z)
 &=\psi(z+1)-\psi(z+1/2)-\frac1{2z}\\
 &=-\int_0^\infty \e^{-2zu}\tanh(u/2)\,\dd u.
\end{align*}
Both \(h(z)\) and the proposed integral tend to zero as
\(\Re z\to+\infty\).  Integrating the derivative from infinity to \(z\),
then replacing \(2z\) by \(x\), gives \cref{eq:D-laplace}.
Alternatively, differentiation under the integral sign verifies the same
differential equation and the same boundary condition.
\end{proof}

Several useful inequalities are immediate.  Since
\(0<\tanh(t/2)<t/2\) for \(t>0\),
\begin{equation}\label{eq:D-simple-bound}
 0<D(x)<\frac1{4x}\qquad(x>0).
\end{equation}
Moreover \(D'(x)<0\), and complete monotonicity follows:
\begin{equation}\label{eq:D-complete-monotone}
 (-1)^mD^{(m)}(x)>0\qquad(m=0,1,2,\ldots;\ x>0).
\end{equation}

\subsection{The complete Borel singularity lattice}

The partial-fraction expansion of the hyperbolic tangent gives
\begin{equation}\label{eq:Borel-partial-fraction}
 \widehat D(t)
 =2\sum_{k=0}^\infty
   \frac1{t^2+\pi^2(2k+1)^2}.
\end{equation}
Indeed both sides are even meromorphic functions, have the same poles and
residues, and agree at the origin.  Thus the Borel singularities are the
simple poles
\begin{equation}\label{eq:Borel-poles}
 t=\pm\ii\pi(2k+1),\qquad k=0,1,2,\ldots,
\end{equation}
with residue \(1/t\) at each pole.  The nearest singularities are at
\(\pm\ii\pi\), so the formal expansion is Gevrey--1 with singulant modulus
\(\pi\).

\begin{remark}[What ``full transseries'' means on the positive axis]
The positive real Laplace ray meets no Borel singularity, so
\cref{eq:D-laplace} is unambiguous and equals the exact function.  The
quantity \(\e^{-\pi x}\) that appears in optimal truncation is the envelope
set by the distance \(\pi\) of the nearest off-axis Borel poles.  It should
not be confused with an additional ambiguous positive-axis sector.  Actual
Stokes jumps arise only when the Laplace direction is rotated through the
imaginary Borel rays; this is treated in \cref{sec:stokes}.
\end{remark}

\section{The logarithmic asymptotic series and exact remainder}

\subsection{General coefficient formula}

We use Bernoulli numbers with convention
\begin{equation}\label{eq:Bernoulli-convention}
 \frac{z}{\e^z-1}=\sum_{n\ge0}B_n\frac{z^n}{n!},
 \qquad B_1=-\frac12.
\end{equation}

\begin{theorem}[All logarithmic coefficients]\label{thm:d-coefficients}
As \(x\to\infty\) in any closed subsector of \(|\arg x|<\pi/2\),
\begin{equation}\label{eq:D-asymptotic}
 D(x)\sim\sum_{r=0}^\infty\frac{d_{2r+1}}{x^{2r+1}},
 \qquad d_{2r}=0,
\end{equation}
where
\begin{align}
 d_{2r+1}
 &=\frac{(2^{2r+2}-1)B_{2r+2}}
        {(2r+1)(2r+2)}\label{eq:d-Bernoulli}\\
 &=(-1)^r\frac{2(2r)!}{\pi^{2r+2}}
       \left(1-2^{-2r-2}\right)\zeta(2r+2).
       \label{eq:d-zeta}
\end{align}
Equivalently, the formal Borel transform is the Taylor expansion
\begin{equation}\label{eq:Borel-Taylor}
 \widehat D(t)=\sum_{r\ge0}d_{2r+1}\frac{t^{2r}}{(2r)!}
 =\frac{\tanh(t/2)}{2t}.
\end{equation}
\end{theorem}

\begin{proof}
Expand each summand in \cref{eq:Borel-partial-fraction} at the origin:
\[
 \frac1{t^2+a^2}
 =\sum_{r=0}^{N-1}\frac{(-1)^rt^{2r}}{a^{2r+2}}
  +\frac{(-1)^Nt^{2N}}{a^{2N}(t^2+a^2)}.
\]
Termwise Laplace integration of the finite sum gives
\[
 d_{2r+1}=2(-1)^r(2r)!
 \sum_{k=0}^\infty\frac1{\pi^{2r+2}(2k+1)^{2r+2}},
\]
which is \cref{eq:d-zeta}.  Euler's formula for \(\zeta(2r+2)\) yields
\cref{eq:d-Bernoulli}.  The same calculation proves the asymptotic
statement, while the unexpanded last term gives the exact remainder below.
\end{proof}

The beginning of the series is
\begin{equation}\label{eq:D-first}
 D(x)\sim
 \frac1{4x}-\frac1{24x^3}+\frac1{20x^5}
 -\frac{17}{112x^7}+\frac{31}{36x^9}
 -\frac{691}{88x^{11}}+\frac{5461}{52x^{13}}-\cdots.
\end{equation}
Substitution in \cref{eq:integer-log-exact} recovers and extends the
logarithmic expansion recorded in OEIS A219931.

\subsection{A signed next-term remainder theorem}

\begin{theorem}[Exact remainder and strict next-term bound]\label{thm:D-remainder}
For \(x>0\) and \(N\ge0\), write
\begin{equation}\label{eq:D-remainder-def}
 D(x)=\sum_{r=0}^{N-1}\frac{d_{2r+1}}{x^{2r+1}}+R_N(x).
\end{equation}
Then, with \(a_k=\pi(2k+1)\),
\begin{equation}\label{eq:D-remainder-integral}
 R_N(x)=2(-1)^N\sum_{k=0}^\infty a_k^{-2N}
 \int_0^\infty\frac{\e^{-xt}t^{2N}}{t^2+a_k^2}\,\dd t,
\end{equation}
and
\begin{equation}\label{eq:D-next-term-bound}
 0<(-1)^NR_N(x)
 <\frac{|d_{2N+1}|}{x^{2N+1}}.
\end{equation}
\end{theorem}

\begin{proof}
Insert the finite geometric decomposition used in the proof of
\cref{thm:d-coefficients} into \cref{eq:Borel-partial-fraction}, then into
the exact Laplace integral.  Every integral in
\cref{eq:D-remainder-integral} is positive.  Replacing
\((t^2+a_k^2)^{-1}\) by the strict upper bound \(a_k^{-2}\), and evaluating
\(\int_0^\infty\e^{-xt}t^{2N}\dd t=(2N)!x^{-2N-1}\), gives exactly the
magnitude of the next term.
\end{proof}

This theorem is stronger than a bare Poincar\'e estimate: every finite
truncation brackets \(D(x)\) with its successor.  It also provides a simple
certificate for all numerical calculations in later sections.

\subsection{Large order and optimal truncation}

From \cref{eq:d-zeta},
\begin{equation}\label{eq:d-large-order}
 d_{2r+1}\sim(-1)^r\frac{2(2r)!}{\pi^{2r+2}}.
\end{equation}
The ratio of successive nonzero term magnitudes is asymptotic to
\((2r+2)(2r+1)/(\pi^2x^2)\).  Hence the least term occurs near
\begin{equation}\label{eq:optimal-r}
 r\sim\frac{\pi x}{2}.
\end{equation}
Stirling's formula then gives the least-term envelope
\begin{equation}\label{eq:least-term-envelope}
 \frac{2\sqrt2}{\pi\sqrt{x}}\e^{-\pi x}
 \left(1+\Ord(x^{-1})\right).
\end{equation}
By \cref{thm:D-remainder}, truncation at or immediately before the least
term attains an error of this scale without any hidden constant explosion.

\section{Forward transseries and Bell-polynomial coefficients}

\subsection{Exponential coefficient calculus}

Let \(Y_n\) denote the complete exponential Bell polynomial, characterized
by
\begin{equation}\label{eq:Bell-egf}
 \exp\left(\sum_{j\ge1}x_j\frac{t^j}{j!}\right)
 =\sum_{n\ge0}Y_n(x_1,\ldots,x_n)\frac{t^n}{n!}.
\end{equation}

\begin{theorem}[All forward coefficients]\label{thm:forward-coeff}
Define \(A_0=1\) and, for \(n\ge1\),
\begin{align}
 A_n
 &=\coeff{t}{n}\exp\left(\sum_{j\ge1}d_jt^j\right)
 \label{eq:A-coeff}\\
 &=\frac1{n!}Y_n(1!d_1,2!d_2,\ldots,n!d_n)\label{eq:A-Bell}\\
 &=\sum_{\substack{k_1,k_2,\ldots\ge0\\
                    \sum_{j\ge1}jk_j=n}}
    \prod_{j\ge1}\frac{d_j^{k_j}}{k_j!}.
    \label{eq:A-partition}
\end{align}
Then
\begin{equation}\label{eq:forward-asymptotic}
 \mathcal S_\sigma(x)
 \sim C_\sigma2^xx^{\sigma/2}
      \sum_{n\ge0}\frac{\sigma^nA_n}{x^n}.
\end{equation}
The coefficients obey the triangular recurrence
\begin{equation}\label{eq:A-recurrence}
 A_0=1,\qquad
 nA_n=\sum_{j=1}^n jd_jA_{n-j}
     =\sum_{\substack{1\le j\le n\\j\text{ odd}}}jd_jA_{n-j}.
\end{equation}
\end{theorem}

\begin{proof}
Exponentiating \cref{eq:D-asymptotic} and using
\cref{eq:Bell-egf} gives \cref{eq:A-Bell}; expanding the exponential into a
product gives the partition sum.  Differentiating
\(A(t)=\exp(\sum d_jt^j)\) yields
\(A'(t)=(\sum jd_jt^{j-1})A(t)\), and coefficient comparison proves
\cref{eq:A-recurrence}.  Finally, because \(D(t)\) is an odd formal series,
\[
 \exp(-D(t))=\exp(D(-t))=\sum_{n\ge0}(-1)^nA_nt^n.
\]
This explains the factor \(\sigma^n\) in \cref{eq:forward-asymptotic}.
\end{proof}

\begin{longtable}{r@{\qquad}l@{\qquad}r@{\qquad}l}
\caption{Universal coefficients \(A_n\) in \(\exp D\).}\label{tab:A}\cr
\toprule
\(n\)&\(A_n\)&\(n\)&\(A_n\)\\
\midrule
\endfirsthead
\toprule
\(n\)&\(A_n\)&\(n\)&\(A_n\)\\
\midrule
\endhead
0&\(1\)&7&\(-39325/262144\)\\
1&\(1/4\)&8&\(-334477/8388608\)\\
2&\(1/32\)&9&\(28717403/33554432\)\\
3&\(-5/128\)&10&\(59697183/268435456\)\\
4&\(-21/2048\)&11&\(-8400372435/1073741824\)\\
5&\(399/8192\)&12&\(-34429291905/17179869184\)\\
6&\(869/65536\)&13&\(7199255611995/68719476736\)\\
\bottomrule
\end{longtable}

\subsection{The two explicit forward sectors}

The even branch is
\begin{align}\label{eq:even-forward}
 \mathcal S_-(x)
 &\sim\sqrt{\frac2\pi}\,2^xx^{-1/2}
 \biggl(
 1-\frac1{4x}+\frac1{32x^2}+\frac5{128x^3}
 -\frac{21}{2048x^4}-\frac{399}{8192x^5}\\[-1mm]
 &\hspace{44mm}
 +\frac{869}{65536x^6}+\frac{39325}{262144x^7}
 -\frac{334477}{8388608x^8}+\cdots
 \biggr).\notag
\end{align}
At \(x=2m\) this is exactly the classical central-binomial expansion in the
variable \(2m\).

The odd branch is
\begin{align}\label{eq:odd-forward}
 \mathcal S_+(x)
 &\sim\frac{2^x x^{1/2}}{\sqrt{2\pi}}
 \biggl(
 1+\frac1{4x}+\frac1{32x^2}-\frac5{128x^3}
 -\frac{21}{2048x^4}+\frac{399}{8192x^5}\\[-1mm]
 &\hspace{44mm}
 +\frac{869}{65536x^6}-\frac{39325}{262144x^7}
 -\frac{334477}{8388608x^8}+\cdots
 \biggr).\notag
\end{align}
For the integer sequence, combine these as
\begin{equation}\label{eq:integer-forward-unified}
 \SF(n)\sim
 \frac{2^n}{\sqrt\pi}\,2^{-\sigma_n/2}n^{\sigma_n/2}
 \sum_{k\ge0}\frac{\sigma_n^kA_k}{n^k},
 \qquad \sigma_n=(-1)^{n+1}.
\end{equation}
This is a periodic or two-phase transseries: the discrete transmonomial
\(\sigma_n\) has magnitude one and must remain visible at every order.

\subsection{Gevrey character and Borel summability}

The exact Borel representation proves that the logarithmic series is
1-summable in every direction that avoids the rays through
\(\pm\ii\pi(2k+1)\).  Exponentiation preserves Gevrey--1 character and
sectorial summability.  Therefore \cref{eq:forward-asymptotic} is not merely
a formal manipulation of Stirling series: in the right half-plane it is the
asymptotic expansion of the exact gamma-ratio branch
\cref{eq:branch-gamma}.  For fixed \(N\),
\begin{equation}\label{eq:forward-error-fixed}
 \mathcal S_\sigma(x)=C_\sigma2^xx^{\sigma/2}
 \left(\sum_{n=0}^{N-1}\frac{\sigma^nA_n}{x^n}
 +\Ord_N(x^{-N})\right).
\end{equation}
After near-least-term truncation, the relative error has the same
\(\e^{-\pi x}/\sqrt{x}\) envelope as the logarithmic correction, although
its hyperasymptotic prefactor is modified by exponentiation.

\section{A universal coefficient calculus for exponential--power inversion}
\label{sec:universal-calculus}

The inverse calculation is clearer if the swing factorial is placed inside a
slightly more general class.  Let \(a>0\), \(b\) be fixed, and suppose
\begin{equation}\label{eq:general-model}
 F(x)=C\e^{ax}x^b\exp Q(x^{-1}),
 \qquad Q(t)=\sum_{j\ge1}q_jt^j.
\end{equation}
At the formal level \(Q\) may be any series over a commutative
\(\mathbb Q\)-algebra.  Analytically it may be a Gevrey asymptotic series
with a chosen sectorial sum.  Taking logarithms gives
\begin{equation}\label{eq:general-log-equation}
 L:=\log(y/C)=ax+b\log x+Q(x^{-1}).
\end{equation}
The swing branches correspond to
\begin{equation}\label{eq:swing-specialization}
 a=\lambda,
 \qquad b=\frac\sigma2,
 \qquad q_j=\sigma d_j.
\end{equation}

\subsection{A parameterized Lagrange lemma}

\begin{lemma}[Lagrange fixed-point formula]\label{lem:fixed-point-LB}
Let \(R\) be a commutative \(\mathbb Q\)-algebra and
\(\Psi(t,u)\in tR[[t,u]]\).  The equation
\begin{equation}\label{eq:fixed-point}
 U(t)=\Psi(t,U(t))
\end{equation}
has a unique solution \(U\in tR[[t]]\), and
\begin{equation}\label{eq:fixed-point-coeff}
 [t^n]U(t)=
 \sum_{m=1}^n\frac1m[t^nu^{m-1}]\Psi(t,u)^m.
\end{equation}
\end{lemma}

\begin{proof}
Introduce an auxiliary marker \(z\) and solve
\(U=z\Psi(t,U)\) in \(R[[t,z]]\).  Ordinary Lagrange--B\"urmann inversion in
\(z\) gives
\[
 U=\sum_{m\ge1}\frac{z^m}{m}[u^{m-1}]\Psi(t,u)^m.
\]
Because \(\Psi\) is divisible by \(t\), only \(m\le n\) contributes to
\([t^n]\); hence setting \(z=1\) is coefficientwise legitimate and gives
\cref{eq:fixed-point-coeff}.  Uniqueness also follows recursively in the
\(t\)-adic topology.
\end{proof}

\subsection{Ordinary partial Bell polynomials}

For a sequence \(\gamma=(\gamma_1,\gamma_2,\ldots)\), define
\begin{equation}\label{eq:ordinary-Bell-def}
 \widehat B_{n,k}(\gamma)
 :=[t^n]\left(\sum_{r\ge1}\gamma_rt^r\right)^k.
\end{equation}
Thus \(\widehat B_{0,0}=1\), \(\widehat B_{n,0}=0\) for \(n>0\), and
\begin{equation}\label{eq:ordinary-Bell-partition}
 \widehat B_{n,k}(\gamma)
 =\sum_{\substack{m_1,m_2,\ldots\ge0\\
                   \sum m_r=k,\ \sum rm_r=n}}
   \frac{k!}{\prod_rm_r!}\prod_r\gamma_r^{m_r}.
\end{equation}
These are ordinary, rather than exponential, partial Bell polynomials.  They
are the natural bookkeeping device for powers of a correction series whose
coefficients are not divided by factorials.

\section{Lambert-normalized inverse transseries}
\label{sec:lambert-inverse}

\subsection{Exact inversion of the dominant block}

Suppress \(Q\) in \cref{eq:general-log-equation}.  The equation
\begin{equation}\label{eq:dominant-equation}
 L=aX+b\log X
\end{equation}
has, when \(b\ne0\), the exact Lambert form
\begin{equation}\label{eq:general-Lambert-core}
 X_k(L)=\frac ba
 W_k\!\left(\frac ab\e^{L/b}\right).
\end{equation}
For the swing parameters this becomes
\begin{equation}\label{eq:swing-Lambert-general}
 X_{\sigma,k}(y)=\frac{\sigma}{2\lambda}
 W_k\!\left(2\sigma\lambda\e^{2\sigma L_\sigma(y)}\right),
 \qquad
 L_\sigma(y)=\log\frac{y}{C_\sigma}.
\end{equation}
The real large branches are
\begin{equation}\label{eq:real-Lambert-cores}
 \boxed{
 X_+(y)=\frac{W_0(4\pi\lambda y^2)}{2\lambda},
 \qquad
 X_-(y)=-\frac{W_{-1}(-4\lambda/(\pi y^2))}{2\lambda}.}
\end{equation}
The appearance of \(W_{-1}\) in the even phase is essential.  Its argument
approaches zero through negative values, and \(W_{-1}(z)\to-\infty\), which
produces the large positive root.  The principal branch \(W_0\) would
produce the small root of the simplified equation and is irrelevant to the
large-index inverse.

\subsection{All constant coefficients around the Lambert core}

\begin{theorem}[Lambert-centered reversion]\label{thm:lambert-centered}
Let \(X\) solve \cref{eq:dominant-equation}, let \(t=X^{-1}\), and seek the
inverse of \cref{eq:general-log-equation} in the form
\begin{equation}\label{eq:Lambert-centered-ansatz}
 x=X+U(t),\qquad U(t)=\sum_{n\ge1}c_nt^n.
\end{equation}
Define
\begin{equation}\label{eq:Psi-general}
 \Psi(t,u):=-\frac1a\left[
 b\log(1+tu)+Q\!\left(\frac{t}{1+tu}\right)
 \right].
\end{equation}
Then \(U=\Psi(t,U)\), and for every \(n\ge1\),
\begin{equation}\label{eq:c-Lagrange-general}
 c_n=\sum_{m=1}^n\frac1m[t^nu^{m-1}]\Psi(t,u)^m.
\end{equation}

Equivalently, put \(\gamma_1=0\) and
\(\gamma_r=c_{r-1}\) for \(r\ge2\).  Then the coefficients are generated
triangularly by
\begin{align}
 c_n=-\frac1a\Bigg[&
 b\sum_{k=1}^n\frac{(-1)^{k+1}}{k}
       \widehat B_{n,k}(\gamma)\notag\\
 &+\sum_{j=1}^nq_j
   \sum_{k=0}^{n-j}(-1)^k\binom{j+k-1}{k}
       \widehat B_{n-j,k}(\gamma)
 \Bigg].\label{eq:c-Bell-recurrence}
\end{align}
Only \(c_1,\ldots,c_{n-1}\) occur on the right-hand side.
\end{theorem}

\begin{proof}
Subtract \(aX+b\log X=L\) from
\(a(X+U)+b\log(X+U)+Q((X+U)^{-1})=L\).  Since
\(X+U=X(1+tU)\), the result is precisely \(U=\Psi(t,U)\).
\Cref{eq:c-Lagrange-general} is
\cref{lem:fixed-point-LB}.  To obtain \cref{eq:c-Bell-recurrence}, write
\(w=tU=\sum_{r\ge1}\gamma_rt^r\), expand
\[
 \log(1+w)=\sum_{k\ge1}\frac{(-1)^{k+1}}{k}w^k,
 \qquad
 (1+w)^{-j}=\sum_{k\ge0}(-1)^k\binom{j+k-1}{k}w^k,
\]
and extract \([t^n]\).
\end{proof}

\begin{remark}
The dominant logarithm has been resummed exactly into \(X\).  Therefore the
remaining \(c_n\) are constants in the coefficient ring; no powers of
\(\log L\) occur.  This is the principal practical advantage over direct
power--logarithmic reversion.
\end{remark}

\subsection{Specialization to the swing factorial}

For \cref{eq:swing-specialization}, define
\begin{equation}\label{eq:Psi-swing}
 \Psi_\sigma(t,u)=-\frac1\lambda\left[
 \frac\sigma2\log(1+tu)
 +\sigma\sum_{j\ge1}d_j
       \left(\frac{t}{1+tu}\right)^j
 \right].
\end{equation}
Then
\begin{equation}\label{eq:c-swing-master}
 c_n^{(\sigma)}=
 \sum_{m=1}^n\frac1m[t^nu^{m-1}]\Psi_\sigma(t,u)^m.
\end{equation}
The first coefficients, reduced with \(\sigma^2=1\), are
\begin{align}
 c_1^{(\sigma)}&=-\frac{\sigma}{4\lambda},
 &c_2^{(\sigma)}&=\frac1{8\lambda^2},\label{eq:c12}\\
 c_3^{(\sigma)}&=
 \frac{2\sigma\lambda^2-3\lambda-3\sigma}{48\lambda^3},
 &c_4^{(\sigma)}&=
 \frac{-4\lambda^2+15\sigma\lambda+6}{192\lambda^4},
 \label{eq:c34}\\
 c_5^{(\sigma)}&=-\frac{
 96\sigma\lambda^4-80\lambda^3+40\sigma\lambda^2
 +135\lambda+30\sigma}{1920\lambda^5},\label{eq:c5}\\
 c_6^{(\sigma)}&=\frac{
 96\lambda^4-180\sigma\lambda^3+215\lambda^2
 +210\sigma\lambda+30}{3840\lambda^6}.
 \label{eq:c6}
\end{align}
Thus the odd inverse begins
\begin{align}\label{eq:odd-inverse-Lambert}
 \mathcal I_+(y)\sim X_+
 &-\frac{1}{4\lambda X_+}
 +\frac{1}{8\lambda^2X_+^2}
 +\frac{2\lambda^2-3\lambda-3}{48\lambda^3X_+^3}\\
 &+\frac{-4\lambda^2+15\lambda+6}{192\lambda^4X_+^4}
 +\cdots,\notag
\end{align}
and the even inverse begins
\begin{align}\label{eq:even-inverse-Lambert}
 \mathcal I_-(y)\sim X_-
 &+\frac{1}{4\lambda X_-}
 +\frac{1}{8\lambda^2X_-^2}
 +\frac{-2\lambda^2-3\lambda+3}{48\lambda^3X_-^3}\\
 &+\frac{-4\lambda^2-15\lambda+6}{192\lambda^4X_-^4}
 +\cdots.\notag
\end{align}
The sign of the first correction follows already from \(D(x)>0\): the odd
phase has \(+D\) in its logarithm and therefore reaches a fixed value at an
index smaller than \(X_+\), while the even phase has \(-D\) and reaches it
at an index larger than \(X_-\).

\subsection{Complex inverse branches}

Every Lambert branch in \cref{eq:swing-Lambert-general} that satisfies
\(|X_{\sigma,k}|\to\infty\) in a sector supplies a formal complex inverse
branch
\begin{equation}\label{eq:complex-inverse-branches}
 \mathcal I_{\sigma,k}(y)
 \sim X_{\sigma,k}(y)+
      \sum_{n\ge1}c_n^{(\sigma)}X_{\sigma,k}(y)^{-n}.
\end{equation}
The coefficient sequence is independent of \(k\); all branch dependence is
carried by \(W_k\) and by the chosen logarithm in \(L_\sigma\).  Analytic
validity requires that the path remain inside a sector avoiding poles of the
gamma quotient and Stokes directions of its Borel transform.  The standard
branch structure and large-argument expansion of \(W_k\) are described in
\cite{CorlessLambertW,DLMFLambertW}.

\section{Direct power--logarithmic inverse transseries}
\label{sec:direct-inverse}

Lambert normalization is compact and numerically efficient, but it hides the
explicit polynomial--logarithmic scale.  We now expand it completely.  For
the general model \cref{eq:general-model}, put
\begin{equation}\label{eq:L-H-general}
 L=\log(y/C),\qquad H=\log(L/a).
\end{equation}
We order monomials lexicographically by powers of \(L^{-1}\), with
\(H=\Ord(\log L)\) treated as a slow variable.

\subsection{Universal coefficient formula}

\begin{theorem}[Direct reversion at infinity]\label{thm:direct-reversion}
The formal inverse of \cref{eq:general-log-equation} has the unique form
\begin{equation}\label{eq:direct-ansatz}
 x\sim\frac1a\left[
 L-bH+\sum_{n\ge1}\frac{P_n(H)}{L^n}
 \right].
\end{equation}
Let
\begin{equation}\label{eq:Phi-direct}
 \Phi(t,u;H):=-b\log\bigl(1+t(u-bH)\bigr)
 -Q\!\left(\frac{at}{1+t(u-bH)}\right).
\end{equation}
Then
\begin{equation}\label{eq:P-Lagrange}
 P_n(H)=\sum_{m=1}^n\frac1m
 [t^nu^{m-1}]\Phi(t,u;H)^m.
\end{equation}

For a triangular Bell recurrence, set
\begin{equation}\label{eq:gamma-direct}
 \gamma_1=-bH,
 \qquad \gamma_r=P_{r-1}(H)\quad(r\ge2).
\end{equation}
Then
\begin{align}
 P_n(H)={}&-b\sum_{k=1}^n\frac{(-1)^{k+1}}{k}
             \widehat B_{n,k}(\gamma)\notag\\
 &-\sum_{j=1}^nq_ja^j
   \sum_{k=0}^{n-j}(-1)^k\binom{j+k-1}{k}
             \widehat B_{n-j,k}(\gamma).
 \label{eq:P-Bell-recurrence}
\end{align}
In particular \(P_n\) has degree at most \(n\).  If \(b\ne0\), its leading
term is
\begin{equation}\label{eq:P-leading}
 P_n(H)=\frac{b^{n+1}}{n}H^n+\text{terms of degree at most }n-1.
\end{equation}
\end{theorem}

\begin{proof}
Set \(t=L^{-1}\) and \(V(t)=\sum_{n\ge1}P_n(H)t^n\).  The ansatz gives
\[
 ax=t^{-1}-bH+V,
 \qquad
 x=\frac1{at}\bigl(1+t(V-bH)\bigr).
\]
Since \(\log(1/(at))=H\), substitution in
\(t^{-1}=ax+b\log x+Q(x^{-1})\) cancels the two occurrences of \(bH\) and
leaves
\[
 V=\Phi(t,V;H).
\]
\Cref{eq:P-Lagrange} follows from
\cref{lem:fixed-point-LB}.  Writing
\(w=t(V-bH)=\sum_{r\ge1}\gamma_rt^r\), then expanding
\(\log(1+w)\) and \((1+w)^{-j}\), gives
\cref{eq:P-Bell-recurrence}.  Induction proves the degree bound.  The
highest power of \(H\) is unaffected by \(Q\); extracting it from the pure
logarithmic equation gives \cref{eq:P-leading}, or follows immediately from
\cref{thm:pure-Lambert-Stirling} below.
\end{proof}

\subsection{Closed Stirling formula for the pure Lambert block}

Let \(\StirlingFirst{n}{m}\) denote the unsigned Stirling number of the first
kind.

\begin{theorem}[Pure Lambert coefficients]\label{thm:pure-Lambert-Stirling}
When \(Q=0\), the polynomials in \cref{eq:direct-ansatz} are
\begin{equation}\label{eq:pure-Lambert-P}
 P_n^{(0)}(H)=
 \frac{b^{n+1}}{n!}
 \sum_{m=1}^n(-1)^{m+1}(m-1)!
 \StirlingFirst{n}{m}\binom{n}{m-1}H^{n-m+1}.
\end{equation}
\end{theorem}

\begin{proof}
With \(Q=0\), \cref{eq:P-Lagrange} becomes
\[
 P_n^{(0)}=
 \sum_{m=1}^n\frac{(-b)^m}{m}
 [t^nu^{m-1}]\log\bigl(1+t(u-bH)\bigr)^m.
\]
Use the signed Stirling expansion
\[
 \log(1+z)^m=m!\sum_{n\ge m}s(n,m)\frac{z^n}{n!},
 \qquad s(n,m)=(-1)^{n-m}\StirlingFirst{n}{m}.
\]
The coefficient of \(t^nu^{m-1}\) is obtained from
\((u-bH)^n\).  Simplifying the signs gives
\cref{eq:pure-Lambert-P}.
\end{proof}

The first pure terms are
\begin{align*}
 P_1^{(0)}&=b^2H,\\
 P_2^{(0)}&=\frac{b^3}{2}(H^2-2H),\\
 P_3^{(0)}&=\frac{b^4}{6}(2H^3-9H^2+6H),\\
 P_4^{(0)}&=\frac{b^5}{12}(3H^4-22H^3+36H^2-12H).
\end{align*}
Formula~\eqref{eq:pure-Lambert-P} is the precise bridge between Lambert
asymptotics and the first-kind Stirling coefficient calculus.

\subsection{Explicit swing inverse polynomials}

For the swing factorial, let
\begin{equation}\label{eq:L-H-swing}
 L_\sigma=\log y+\frac12\log\pi+\frac\sigma2\lambda,
 \qquad H_\sigma=\log\frac{L_\sigma}{\lambda}.
\end{equation}
Then
\begin{equation}\label{eq:direct-swing-inverse}
 \boxed{\quad
 \mathcal I_\sigma(y)\sim\frac1\lambda\left[
 L_\sigma-\frac\sigma2H_\sigma
 +\sum_{n\ge1}\frac{P_n^{(\sigma)}(H_\sigma)}{L_\sigma^n}
 \right].\quad}
\end{equation}
The first five coefficient polynomials are
\begin{align}
 P_1^{(\sigma)}(H)
 &=\frac{H-\sigma\lambda}{4},\label{eq:P1-swing}\\
 P_2^{(\sigma)}(H)
 &=\frac{\sigma H^2-2(\lambda+\sigma)H+2\lambda}{16},
 \label{eq:P2-swing}\\
 P_3^{(\sigma)}(H)
 &=\frac1{96}\Bigl[
 2H^3-(6\sigma\lambda+9)H^2+(18\sigma\lambda+6)H\notag\\
 &\hspace{29mm}
 +4\sigma\lambda^3-6\lambda^2-6\sigma\lambda
 \Bigr],\label{eq:P3-swing}\\
 P_4^{(\sigma)}(H)
 &=\frac1{384}\Bigl[
 3\sigma H^4-(12\lambda+22\sigma)H^3
 +(66\lambda+36\sigma)H^2\notag\\
 &\hspace{15mm}
 +(24\lambda^3-36\sigma\lambda^2-72\lambda-12\sigma)H\notag\\
 &\hspace{29mm}
 -8\lambda^3+30\sigma\lambda^2+12\lambda
 \Bigr],\label{eq:P4-swing}\\
 P_5^{(\sigma)}(H)
 &=\frac1{3840}\Bigl[
 12H^5-(60\sigma\lambda+125)H^4
 +(500\sigma\lambda+350)H^3\notag\\
 &\hspace{13mm}
 +(240\sigma\lambda^3-360\lambda^2-1050\sigma\lambda-300)H^2\notag\\
 &\hspace{13mm}
 +(-280\sigma\lambda^3+780\lambda^2+600\sigma\lambda+60)H\notag\\
 &\hspace{13mm}
 -192\sigma\lambda^5+160\lambda^4-80\sigma\lambda^3
 -270\lambda^2-60\sigma\lambda
 \Bigr].\label{eq:P5-swing}
\end{align}
Every coefficient is obtained by the single recurrence
\cref{eq:P-Bell-recurrence} with \(b=\sigma/2\) and
\(q_j=\sigma d_j\); no order-specific ansatz is needed.

\subsection{Relation between the two inverse normalizations}

The Lambert-centered expansion \cref{eq:Lambert-centered-ansatz} resums the
entire pure block \cref{eq:pure-Lambert-P}.  Expanding \(W_0\) or
\(W_{-1}\) in iterated logarithms and then expanding each power
\(X^{-n}\) reproduces \cref{eq:direct-swing-inverse}.  Conversely, grouping
in the direct series all terms generated solely by
\((\sigma/2)\log x\) reconstructs \(X_\sigma\).

Formally, the direct coefficients live in
\[
 \mathbb Q[\lambda,\lambda^{-1},\sigma,H]/(\sigma^2-1),
\]
whereas the Lambert-centered coefficients live in the smaller constant ring
\[
 \mathbb Q[\lambda,\lambda^{-1},\sigma]/(\sigma^2-1).
\]
This algebraic compression explains the substantially better numerical
behavior observed in \cref{sec:numerics}.

\section{Perturbative Lagrange operators and exact inverse transport}
\label{sec:operator-inversion}

There is a third all-orders description that is especially convenient for
Borel and Stokes analysis.

\begin{theorem}[Inverse perturbation operator]\label{thm:inverse-operator}
Let \(F_0\) be locally invertible, let \(X=F_0^{-1}(L)\), and consider
\begin{equation}\label{eq:perturbed-inverse-problem}
 L=F_0(x)+\varepsilon R(x).
\end{equation}
Define the derivative with respect to the leading coordinate by
\begin{equation}\label{eq:D0-operator}
 \mathscr D_0:=\frac1{F_0'(X)}\frac{\dd}{\dd X}.
\end{equation}
Then, as a formal power series in \(\varepsilon\),
\begin{equation}\label{eq:inverse-operator-formula}
 x=X+\sum_{m\ge1}\frac{(-\varepsilon)^m}{m!}
 \mathscr D_0^{m-1}
 \left[\frac{R(X)^m}{F_0'(X)}\right].
\end{equation}
\end{theorem}

\begin{proof}
Use \(u=F_0(x)\).  \Cref{eq:perturbed-inverse-problem} becomes
\(L=u+\varepsilon r(u)\), where
\(r(u)=R(F_0^{-1}(u))\).  Lagrange--B\"urmann inversion for the composite
function \(F_0^{-1}(u)\) gives
\[
 F_0^{-1}(u(L))=F_0^{-1}(L)
 +\sum_{m\ge1}\frac{(-\varepsilon)^m}{m!}
 \left(\frac{\dd}{\dd L}\right)^{m-1}
 \left[(F_0^{-1})'(L)r(L)^m\right].
\]
Now \((F_0^{-1})'(L)=1/F_0'(X)\) and
\(\dd/\dd L=\mathscr D_0\).
\end{proof}

For the swing problem,
\begin{equation}\label{eq:swing-F0-R}
 F_{0,\sigma}(x)=\lambda x+\frac\sigma2\log x,
 \qquad R_\sigma(x)=\sigma D(x),
 \qquad F_{0,\sigma}'(x)=\lambda+\frac\sigma{2x}.
\end{equation}
Thus, with \(X=X_\sigma(y)\),
\begin{equation}\label{eq:swing-operator-series}
 \mathcal I_\sigma(y)
 =X+\sum_{m\ge1}\frac{(-\sigma)^m}{m!}
 \mathscr D_{0,\sigma}^{m-1}
 \left[\frac{D(X)^m}{\lambda+\sigma/(2X)}\right],
\end{equation}
where \(D\) may be interpreted either as its formal asymptotic series or as
the exact Borel sum \cref{eq:D-laplace}.  The first three perturbative layers
are
\begin{align}
 \mathcal I_\sigma(y)=X
 &-\sigma\frac{D(X)}{F_{0,\sigma}'(X)}
 +\frac12\mathscr D_{0,\sigma}
   \left[\frac{D(X)^2}{F_{0,\sigma}'(X)}\right]\notag\\
 &-\frac\sigma6\mathscr D_{0,\sigma}^2
   \left[\frac{D(X)^3}{F_{0,\sigma}'(X)}\right]+\cdots.
 \label{eq:swing-operator-first}
\end{align}
Expanding \(D(X)\) in inverse powers of \(X\) reproduces exactly the
coefficients \(c_n^{(\sigma)}\).

\section{Borel singularities, Stokes jumps, and inverse transseries}
\label{sec:stokes}

For a direction \(\theta\), write
\begin{equation}\label{eq:lateral-sums}
 \mathcal S_{\theta^\pm}D(x)
 :=\int_{0}^{\e^{\ii(\theta\pm0)}\infty}
   \e^{-xt}\widehat D(t)\,\dd t.
\end{equation}
whenever the rotated Laplace integral converges.  Define the discontinuity
by lower minus upper lateral summation.  The residue theorem and
\cref{eq:Borel-poles} give the exact bridge equations
\begin{align}
 \Disc_{\pi/2}D(x)
 &:=\mathcal S_{(\pi/2)^-}D-
     \mathcal S_{(\pi/2)^+}D
 =2\sum_{k\ge0}\frac{\e^{-\ii\pi(2k+1)x}}{2k+1},
 \label{eq:upper-Stokes}\\
 \Disc_{-\pi/2}D(x)
 &:=\mathcal S_{(-\pi/2)^-}D-
     \mathcal S_{(-\pi/2)^+}D
 =-2\sum_{k\ge0}\frac{\e^{\ii\pi(2k+1)x}}{2k+1}.
 \label{eq:lower-Stokes}
\end{align}
In their respective convergence half-planes these may be summed as
\begin{equation}\label{eq:Stokes-artanh}
 \Disc_{\pi/2}D=2\operatorname{artanh}(\e^{-\ii\pi x}),
 \qquad
 \Disc_{-\pi/2}D=-2\operatorname{artanh}(\e^{\ii\pi x}).
\end{equation}
The logarithm of the swing branch has jump
\begin{equation}\label{eq:logS-jump}
 \Disc_\theta\log\mathcal S_\sigma(x)
 =\sigma\Disc_\theta D(x),
\end{equation}
so the function itself changes multiplicatively by the exponential of this
quantity.

\begin{proposition}[Transport of a Stokes jump through inversion]
\label{prop:inverse-Stokes}
Let \(F(x)\) be one lateral logarithmic branch and let the other be
\(F(x)+\Delta(x)\).  If \(x=F^{-1}(L)\), the inverse on the second side has
the formal expansion
\begin{equation}\label{eq:inverse-Stokes-full}
 x_{\mathrm{new}}=x+
 \sum_{m\ge1}\frac{(-1)^m}{m!}
 \left(\frac1{F'(x)}\frac{\dd}{\dd x}\right)^{m-1}
 \left[\frac{\Delta(x)^m}{F'(x)}\right].
\end{equation}
In particular,
\begin{equation}\label{eq:inverse-Stokes-leading}
 x_{\mathrm{new}}-x
 =-\frac{\Delta(x)}{F'(x)}+	ext{higher exponential sectors}.
\end{equation}
For the swing inverse, take
\(\Delta=\sigma\Disc_\theta D\) and
\(F'=\lambda+\sigma/(2x)+\sigma D'(x)\).
\end{proposition}

\begin{proof}
Apply \cref{thm:inverse-operator} with \(F_0=F\),
\(\varepsilon=1\), and \(R=\Delta\).
\end{proof}

Equations \cref{eq:upper-Stokes,eq:lower-Stokes,eq:inverse-Stokes-full}
complete the resurgent transseries data: the Borel singularity locations,
residues, all lateral jumps, and their nonlinear transport through the
inverse are explicit.  Along the positive real axis the preferred sum is
\cref{eq:D-laplace}; no lateral choice is necessary.

\section{Error control, computation, and numerical validation}
\label{sec:numerics}

\subsection{Residual certification}

Asymptotic inversion should be checked in the original equation, preferably
in logarithmic form to avoid overflow.

\begin{proposition}[Residual-to-index error bound]\label{prop:residual-bound}
Let \(x_*\) approximate \(\mathcal I_\sigma(y)\), and define
\begin{equation}\label{eq:residual}
 r_*:=\log\mathcal S_\sigma(x_*)-\log y.
\end{equation}
If \((\log\mathcal S_\sigma)'(x)\ge m>0\) on the interval joining \(x_*\)
and \(\mathcal I_\sigma(y)\), then
\begin{equation}\label{eq:residual-bound}
 |x_*-\mathcal I_\sigma(y)|\le\frac{|r_*|}{m}.
\end{equation}
\end{proposition}

\begin{proof}
Apply the mean-value theorem to \(\log\mathcal S_\sigma\).
\end{proof}

The exact Laplace representation gives
\begin{equation}\label{eq:Dprime-bound}
 0<-D'(x)<\frac1{4x^2},
\end{equation}
so for \(x\ge1\),
\begin{align}
 (\log\mathcal S_+)'(x)
 &>\lambda+\frac1{2x}-\frac1{4x^2}>\lambda,
 \label{eq:odd-derivative-lower}\\
 (\log\mathcal S_-)'(x)
 &>\lambda-\frac1{2x}.
 \label{eq:even-derivative-lower}
\end{align}
These elementary bounds are sufficient for certified large-index inversion;
tighter bounds are available directly from the digamma expression
\cref{eq:logder}.

\subsection{Newton refinement}

Starting from either asymptotic approximation, Newton's method applied to the
logarithmic equation is
\begin{equation}\label{eq:Newton}
 x_{j+1}=x_j-
 \frac{\log\mathcal S_\sigma(x_j)-\log y}
 {\lambda+\dfrac\sigma2
  [\psi(x_j/2+1)-\psi(x_j/2+1/2)]}.
\end{equation}
Evaluation through \(\log\Gamma\) and \(\psi\) is stable even when
\(\mathcal S_\sigma(x)\) itself is far outside floating-point range.  Once
the asymptotic seed is in its natural regime, one or two Newton steps usually
recover full working precision.

\subsection{Forward accuracy}

Table \ref{tab:forward-errors} uses \(x=20\).  The truncation labeled \(N\)
retains \(A_0,\ldots,A_{N-1}\) in \cref{eq:forward-asymptotic}.

\begin{table}[htbp]
\centering
\caption{Relative errors of the forward expansion at \(x=20\).}
\label{tab:forward-errors}
\begin{tabular}{rcc}
\toprule
\(N\)&even branch \(\sigma=-1\)&odd branch \(\sigma=+1\)\\
\midrule
1  &\(1.2573193\times10^{-2}\)&\(1.2417071\times10^{-2}\)\\
2  &\(8.3971506\times10^{-5}\)&\(7.2284564\times10^{-5}\)\\
4  &\(7.9979454\times10^{-8}\)&\(4.8170528\times10^{-8}\)\\
6  &\(3.2525110\times10^{-10}\)&\(8.8967744\times10^{-11}\)\\
8  &\(3.2108180\times10^{-12}\)&\(9.7020566\times10^{-14}\)\\
10 &\(5.8884939\times10^{-14}\)&\(1.5606307\times10^{-14}\)\\
12 &\(1.7866041\times10^{-15}\)&\(6.7178815\times10^{-16}\)\\
\bottomrule
\end{tabular}
\end{table}

The decrease is far faster than a fixed-order theorem alone suggests because
\(x=20\) is still well before the least-term index, which is near
\(\pi x/2\approx31\) in the nonzero logarithmic series.

\subsection{Inverse accuracy and the advantage of Lambert centering}

For each phase, set \(y=\mathcal S_\sigma(50)\), so that the exact inverse is
\(50\).  In Table \ref{tab:inverse-errors}, \(N=0\) denotes only the leading
core: \((L_\sigma-\sigma H_\sigma/2)/\lambda\) in the direct column and
\(X_\sigma\) in the Lambert column.  For \(N>0\), the first \(N\) correction
terms are retained.

\begin{table}[htbp]
\centering
\caption{Absolute inverse errors for an exact index \(x=50\).}
\label{tab:inverse-errors}
\begin{tabular}{rcccc}
\toprule
&\multicolumn{2}{c}{even phase \(\sigma=-1\)}
&\multicolumn{2}{c}{odd phase \(\sigma=+1\)}\\
\cmidrule(lr){2-3}\cmidrule(lr){4-5}
\(N\)&direct&Lambert&direct&Lambert\\
\midrule
0  &\(4.92291\times10^{-2}\)&\(7.31859\times10^{-3}\)
   &\(3.24899\times10^{-2}\)&\(7.11042\times10^{-3}\)\\
1  &\(9.28054\times10^{-4}\)&\(1.04056\times10^{-4}\)
   &\(2.43401\times10^{-4}\)&\(1.02029\times10^{-4}\)\\
2  &\(8.18595\times10^{-6}\)&\(4.29703\times10^{-8}\)
   &\(4.74681\times10^{-6}\)&\(2.00947\times10^{-6}\)\\
3  &\(6.38122\times10^{-7}\)&\(2.27690\times10^{-8}\)
   &\(2.67151\times10^{-7}\)&\(5.08233\times10^{-8}\)\\
4  &\(4.81875\times10^{-8}\)&\(5.65588\times10^{-11}\)
   &\(4.49650\times10^{-9}\)&\(1.40431\times10^{-9}\)\\
6  &\(8.92594\times10^{-12}\)&\(2.85289\times10^{-13}\)
   &\(9.28838\times10^{-12}\)&\(8.11044\times10^{-13}\)\\
8  &\(2.71054\times10^{-13}\)&\(8.38506\times10^{-16}\)
   &\(5.17449\times10^{-15}\)&\(1.18162\times10^{-15}\)\\
10 &\(1.56412\times10^{-16}\)&\(3.13520\times10^{-18}\)
   &\(1.42623\times10^{-17}\)&\(8.47275\times10^{-19}\)\\
\bottomrule
\end{tabular}
\end{table}

The nonmonotone error at a few low truncation orders is normal for an
asymptotic series.  The overall comparison confirms the structural point:
Lambert centering has already resummed the largest tower of
\((\log L)^m/L^n\) terms.

\subsection{Recovering integer indices}

For an approximate equation \(\SF(n)\approx y\), compute both continuous
inverse branches and project to their parity lattices:
\begin{equation}\label{eq:nearest-parity}
 n_-^{\mathrm{near}}=2\left\lfloor\frac{\mathcal I_-(y)}2+\frac12\right\rfloor,
 \qquad
 n_+^{\mathrm{near}}=2\left\lfloor
       \frac{\mathcal I_+(y)-1}{2}+\frac12\right\rfloor+1.
\end{equation}
Evaluate \(\log\SF(n)\) at these candidates and their nearest same-parity
neighbors, then choose according to the desired criterion.  For a threshold
query, the greatest same-parity indices satisfying \(\SF(n)\le y\) are
obtained from
\begin{equation}\label{eq:floor-parity}
 n_-^{\le}=2\left\lfloor\frac{\mathcal I_-(y)}2\right\rfloor,
 \qquad
 n_+^{\le}=2\left\lfloor\frac{\mathcal I_+(y)-1}{2}\right\rfloor+1,
\end{equation}
with a one-step exact correction if roundoff places an approximation across
a boundary.

\section{Further consequences of the two-phase inverse}

\subsection{Separation of the two inverse branches}

For a common value \(y\), define
\begin{equation}\label{eq:Y-h-common}
 Y:=\log(y\sqrt\pi),\qquad h:=\log(Y/\lambda).
\end{equation}
Then \(L_\sigma=Y+\sigma\lambda/2\).  Subtracting the first terms of
\cref{eq:direct-swing-inverse} gives
\begin{equation}\label{eq:branch-gap}
 \boxed{\quad
 \mathcal I_-(y)-\mathcal I_+(y)
 =\frac{h}{\lambda}-1+\frac1{2Y}
 +\Ord\!\left(\frac{h^2}{Y^2}\right).
 \quad}
\end{equation}
In particular,
\begin{equation}\label{eq:branch-gap-leading}
 \mathcal I_-(y)-\mathcal I_+(y)
 =\log_2\!\left(\frac{\log(y\sqrt\pi)}{\log2}\right)-1+o(1).
\end{equation}
Thus the even index producing a given large magnitude exceeds the odd index
by a slowly growing \(\log\log y\) distance.  This is the inverse reflection
of the factor of order \(x/2\) between the two forward branches at a common
argument.

\subsection{The counting function}

Let
\begin{equation}\label{eq:counting-def}
 N(y):=\#\{n\in\mathbb N_0:\ \SF(n)\le y\}.
\end{equation}
Each parity subsequence is monotone, so
\begin{equation}\label{eq:counting-inverses}
 N(y)=\frac{\mathcal I_-(y)+\mathcal I_+(y)}2+\Ord(1).
\end{equation}
The \(\pm(1/2)\log x\) contributions cancel to leading order in the average,
and therefore
\begin{equation}\label{eq:counting-leading}
 \boxed{\qquad
 N(y)=\frac{\log(y\sqrt\pi)}{\log2}+\Ord(1)
 =\log_2(y\sqrt\pi)+\Ord(1).
 \qquad}
\end{equation}
A smooth version of the counting function can be expanded to arbitrary order
by averaging the two polynomial families
\(P_n^{(-)}\) and \(P_n^{(+)}\) after substituting
\(L_\sigma=Y+\sigma\lambda/2\).

\section{Symbolic coefficient generation}
\label{sec:symbolic}

The following SymPy program implements the formulae of the article.  It
constructs \(d_j\), the forward Bell coefficients \(A_n\), the
Lambert-centered coefficients \(c_n^{(\sigma)}\), and the direct
polynomials \(P_n^{(\sigma)}(H)\).  Reduction modulo \(\sigma^2-1\) keeps
the two parity phases in a single expression.

\begin{lstlisting}[style=pythonstyle,caption={All-order coefficient generation.}]
import sympy as sp

N = 12
t, u, H = sp.symbols("t u H")
lam, sigma = sp.symbols("lambda sigma", nonzero=True)

# Logarithmic coefficients d_j; absent even coefficients are zero.
d = {
    2*r + 1:
        (2**(2*r + 2) - 1) * sp.bernoulli(2*r + 2)
        / ((2*r + 1) * (2*r + 2))
    for r in range(N)
}

D = sum(d.get(j, 0) * t**j for j in range(1, N + 1))
E = sp.series(sp.exp(D), t, 0, N + 1).removeO().expand()
A = [sp.factor(E.coeff(t, n)) for n in range(N + 1)]


def parity_reduce(expr):
    """Reduce a polynomial expression modulo sigma**2 - 1."""
    poly = sp.Poly(sp.expand(expr), sigma)
    ans = 0
    for (power,), coefficient in poly.terms():
        ans += coefficient * (sigma if power % 2 else 1)
    return sp.factor(ans)


# Lambert-centered inverse: U = Psi(t,U), U = sum c_n t^n.
U = sp.Integer(0)
c = []
for n in range(1, N + 1):
    Q_at_shift = sigma * sum(
        d.get(j, 0) * (t / (1 + t*U))**j
        for j in range(1, N + 1)
    )
    Psi = -(sigma*sp.log(1 + t*U)/2 + Q_at_shift) / lam
    cn = sp.series(Psi, t, 0, n + 1).removeO().expand().coeff(t, n)
    cn = parity_reduce(cn)
    c.append(cn)
    U += cn * t**n


# Direct inverse: V = Phi(t,V;H), V = sum P_n(H)t^n.
V = sp.Integer(0)
P = []
for n in range(1, N + 1):
    w = t * (V - sigma*H/2)
    Q_at_direct_shift = sigma * sum(
        d.get(j, 0) * (lam*t / (1 + w))**j
        for j in range(1, N + 1)
    )
    Phi = -sigma*sp.log(1 + w)/2 - Q_at_direct_shift
    pn = sp.series(Phi, t, 0, n + 1).removeO().expand().coeff(t, n)
    pn = parity_reduce(pn)
    P.append(pn)
    V += pn * t**n

for n, value in enumerate(A):
    print("A", n, "=", value)
for n, value in enumerate(c, 1):
    print("c", n, "=", value)
for n, value in enumerate(P, 1):
    print("P", n, "=", value)
\end{lstlisting}

The recurrence form is usually faster than direct multivariate coefficient
extraction.  At order \(N\), truncated-series arithmetic gives a
straightforward \(\widetilde{\Ord}(N^2)\) implementation; Newton iteration
for formal power series can reduce this further when thousands of terms are
required.

\section{Conclusions}

The swing factorial is a compact example in which several asymptotic
structures that are often studied separately occur at once.

\begin{enumerate}
\item Its parity split is leading-order information.  The correct forward
      object is a two-component periodic transseries, and the correct inverse
      is a pair of monotone branches.
\item Duplication reduces both branches to one gamma ratio.  After removal of
      its square-root scale, the entire logarithmic correction has the exact
      Borel transform \(\tanh(t/2)/(2t)\).
\item Bernoulli numbers give the logarithmic coefficients; complete Bell
      polynomials exponentiate them; ordinary Bell polynomials organize
      nonlinear reversion; and unsigned first-kind Stirling numbers give the
      pure Lambert inverse block in closed form.
\item The odd real inverse is normalized by \(W_0\), whereas the even real
      inverse is normalized by \(W_{-1}\).  Missing this branch distinction
      selects the wrong root before any asymptotic correction is computed.
\item The Borel pole lattice supplies exact remainder signs, optimal
      truncation, lateral Stokes jumps, and, through the inverse perturbation
      operator, the full exponential-sector transport of every inverse
      branch.
\end{enumerate}

All coefficients in both forward and inverse expansions are therefore
computable by explicit finite formulae at arbitrary order.  No unspecified
``successive matching'' step remains: each layer is either a Bernoulli
number, a Bell polynomial, a Stirling number, or a coefficient extraction
from a stated formal fixed-point equation.

\appendix

\section{Additional Lambert-centered coefficients}

For completeness, the next four coefficients after
\cref{eq:c5,eq:c6} are
\begin{align}
 c_7^{(\sigma)}={}&\frac1{53760\lambda^7}\Bigl(
 8160\sigma\lambda^6-4312\lambda^5+1428\sigma\lambda^4
 +1050\lambda^3\notag\\
 &\hspace{33mm}
 -4025\sigma\lambda^2-2100\lambda-210\sigma\Bigr),
 \label{eq:c7}\\
 c_8^{(\sigma)}={}&\frac1{645120\lambda^8}\Bigl(
 -48960\lambda^6+56056\sigma\lambda^5-40908\lambda^4
 +15015\sigma\lambda^3\notag\\
 &\hspace{33mm}
 +50610\lambda^2+17010\sigma\lambda+1260\Bigr),
 \label{eq:c8}\\
 c_9^{(\sigma)}={}&-\frac1{1290240\lambda^9}\Bigl(
 1111040\sigma\lambda^8-413184\lambda^7+79616\sigma\lambda^6
 +43512\lambda^5\notag\\
 &\hspace{20mm}
 -83748\sigma\lambda^4+83475\lambda^3+92610\sigma\lambda^2
 +22050\lambda+1260\sigma\Bigr),
 \label{eq:c9}\\
 c_{10}^{(\sigma)}={}&\frac1{5160960\lambda^{10}}\Bigl(
 2222080\lambda^8-1756032\sigma\lambda^7+779152\lambda^6
 -203280\sigma\lambda^5\notag\\
 &\hspace{17mm}
 -162582\lambda^4+487305\sigma\lambda^3+309960\lambda^2
 +55440\sigma\lambda+2520\Bigr).
 \label{eq:c10}
\end{align}

\section{A compact derivation from the general gamma-ratio expansion}

The classical Tricomi--Erd\'elyi expansion for a ratio of gamma functions
\cite{TricomiErdelyi,DLMFGamma} implies
\begin{equation}\label{eq:general-gamma-ratio-log}
 \log\frac{\Gamma(z+a)}{\Gamma(z+b)}
 \sim(a-b)\log z+
 \sum_{n\ge1}\frac{(-1)^{n+1}}
 {n(n+1)z^n}\bigl(B_{n+1}(a)-B_{n+1}(b)\bigr),
\end{equation}
where \(B_n(x)\) are Bernoulli polynomials.  Taking
\(a=1\), \(b=1/2\), and \(z=x/2\), then using
\[
 B_m(1/2)=(2^{1-m}-1)B_m,
 \qquad B_{2r+1}=0\quad(r\ge1),
\]
reduces \cref{eq:general-gamma-ratio-log} to
\cref{eq:d-Bernoulli}.  The Laplace proof in the main text is stronger for
this special ratio because it simultaneously supplies the exact Borel
transform and remainder sign.

\begin{thebibliography}{99}

\bibitem{OEISA056040}
The OEIS Foundation,
\emph{A056040: Swinging factorial},
\url{https://oeis.org/A056040}.

\bibitem{OEISA219931}
P. Luschny and contributors,
\emph{A219931: Coefficients related to an asymptotic expansion of the
logarithm of the central binomial},
The On-Line Encyclopedia of Integer Sequences,
\url{https://oeis.org/A219931}.

\bibitem{LuschnySwing}
P. Luschny,
\emph{Die schwingende Fakult\"at und Orbitalsysteme}, 2011,
\url{https://oeis.org/A180000/a180000.pdf}.

\bibitem{TricomiErdelyi}
F. G. Tricomi and A. Erd\'elyi,
The asymptotic expansion of a ratio of gamma functions,
\emph{Pacific Journal of Mathematics} \textbf{1} (1951), 133--142,
\url{https://doi.org/10.2140/pjm.1951.1.133}.

\bibitem{DLMFGamma}
NIST Digital Library of Mathematical Functions,
\S5.11, \emph{Asymptotic Expansions of the Gamma Function},
\url{https://dlmf.nist.gov/5.11}.

\bibitem{CorlessLambertW}
R. M. Corless, G. H. Gonnet, D. E. G. Hare, D. J. Jeffrey, and D. E. Knuth,
On the Lambert \(W\) function,
\emph{Advances in Computational Mathematics} \textbf{5} (1996), 329--359,
\url{https://doi.org/10.1007/BF02124750}.

\bibitem{DLMFLambertW}
NIST Digital Library of Mathematical Functions,
\S4.13, \emph{Lambert \(W\)-Function},
\url{https://dlmf.nist.gov/4.13}.

\bibitem{Comtet}
L. Comtet,
\emph{Advanced Combinatorics},
D. Reidel, Dordrecht, 1974.

\bibitem{FlajoletSedgewick}
P. Flajolet and R. Sedgewick,
\emph{Analytic Combinatorics},
Cambridge University Press, 2009.

\bibitem{Olver}
F. W. J. Olver,
\emph{Asymptotics and Special Functions},
A K Peters, 1997 reprint.

\bibitem{ParisKaminski}
R. B. Paris and D. Kaminski,
\emph{Asymptotics and Mellin--Barnes Integrals},
Cambridge University Press, 2001.

\bibitem{DLMFLagrange}
NIST Digital Library of Mathematical Functions,
\S1.10(vii), \emph{Lagrange Inversion Theorem},
\url{https://dlmf.nist.gov/1.10}.

\bibitem{ProveItCoeff}
V. Reshetnikov,
\emph{Combinatorial Coefficient Calculus}, ProveIt research draft,
\url{https://github.com/VladimirReshetnikov/ProveIt/tree/main/Analysis/FabiusFunction/docs/semi-formalized-research-frontiers/drafts/combinatorial-coefficient-calculus/Combinatorial_Coefficient_Calculus}.

\bibitem{ProveItTransseries}
V. Reshetnikov,
\emph{Series and Transseries Research Drafts}, ProveIt,
\url{https://github.com/VladimirReshetnikov/ProveIt/tree/main/Analysis/FabiusFunction/docs/semi-formalized-research-frontiers/drafts/series-and-transseries}.

\end{thebibliography}

\end{document}
